\chapter{Análisis de sensibilidad}\label{cap:sensibilidad}

Varias decisiones del análisis son razonables pero discutibles: estimar por MCO y no por
mínimos cuadrados ponderados, seleccionar variables, analizar casos completos, conservar los
atípicos o tratar la dependencia por estado sólo a través de los errores típicos. Para comprobar
que las conclusiones no dependen de ellas, el modelo final se reestimó bajo
\res{n-especificaciones} especificaciones alternativas y con un bootstrap por conglomerados
(\cref{tab:sensibilidad,fig:sensibilidad}):

\begin{description}[style=nextline,leftmargin=1.2em]
  \item[Ponderación (MCPF).] Pesos $1/\hat\sigma^2(n)$ con la función de varianza estimada: el
  estimador eficiente si la heterocedasticidad se debe al tamaño del condado.
  \item[Modelo máximo.] Todas las candidatas tras la poda, sin selección de variables
  (\cref{sec:comparacion}).
  \item[Efectos fijos de estado.] Una indicadora por estado: los coeficientes se estiman sólo con la
  variación entre condados del mismo estado, lo que elimina cualquier factor estatal no medido
  (registro, políticas sanitarias, clima). La región desaparece, absorbida por los estados.
  \item[Modelo mixto.] Intercepto aleatorio por estado, que modela explícitamente la dependencia
  intraestatal en lugar de corregirla en los errores típicos.
  \item[Regresión robusta de Huber.] M-estimador que pondera a la baja los residuos grandes
  \parencite{huber1964}: resistente a los atípicos en la respuesta.
  \item[Sin observaciones influyentes.] Reajuste sin los condados con distancia de Cook mayor que
  $4/n$ (\res{n-sin-influyentes} condados restantes).
  \item[Imputación múltiple.] Los \res{n-imputacion} condados, con los ausentes imputados en
  \res{m-imputaciones} conjuntos y los resultados combinados por las reglas de Rubin (\cref{cap:ausentes}).
  \item[Bootstrap por conglomerados.] \res{boot-reps} remuestreos de estados completos con
  reemplazo; los intervalos percentiles no suponen normalidad ni una forma concreta de la
  heterocedasticidad \parencite{efron1993,cameron2015}.
\end{description}

\begin{table}[htbp]
\centering
\footnotesize
\caption{Análisis de sensibilidad: coeficientes de las seis explicativas con mayor $|t|$ bajo especificaciones alternativas}\label{tab:sensibilidad}
\begin{tabularx}{\linewidth}{X r r r r r r r }
\toprule
Especificación & $n$ & \rotatebox{70}{Incidencia de cáncer} & \rotatebox{70}{Edad mediana} & \rotatebox{70}{Solo seguro público} & \rotatebox{70}{Grado universitario (25+)} & \rotatebox{70}{Tamaño medio del hogar} & \rotatebox{70}{Bachillerato (18–24)} \\
\midrule
MCO con errores típicos clásicos (salida tipo SPSS) & \num{2840} & \num{0.187} & \num{-0.788} & \num{0.892} & \num{-1.054} & \num{-12.875} & \num{0.241} \\
MCO con errores típicos HC3 & \num{2840} & \num{0.187} & \num{-0.788} & \num{0.892} & \num{-1.054} & \num{-12.875} & \num{0.241} \\
Principal: MCO + errores cluster por estado & \num{2840} & \num{0.187} & \num{-0.788} & \num{0.892} & \num{-1.054} & \num{-12.875} & \num{0.241} \\
MCPF con $\sigma^2$(n) = a + b/n y errores cluster & \num{2840} & \num{0.178} & \num{-0.929} & \num{0.812} & \num{-1.134} & \num{-16.176} & \num{0.271} \\
Efectos fijos de estado (MCPF + cluster) & \num{2840} & \num{0.188} & \num{-0.581} & \num{0.915} & \num{-1.015} & \num{-10.647} & \num{0.154} \\
Modelo mixto (intercepto aleatorio por estado) & \num{2840} & \num{0.188} & \num{-0.595} & \num{0.856} & \num{-1.118} & \num{-10.680} & \num{0.163} \\
Regresión robusta de Huber (M-estimador) & \num{2840} & \num{0.188} & \num{-0.875} & \num{0.777} & \num{-1.192} & \num{-14.634} & \num{0.252} \\
Sin los 185 condados influyentes (Cook > 4/n) & \num{2655} & \num{0.175} & \num{-0.927} & \num{0.783} & \num{-1.283} & \num{-14.998} & \num{0.256} \\
Imputación múltiple (20 conjuntos, todos los condados) & \num{3047} & \num{0.188} & \num{-0.761} & \num{0.814} & \num{-1.090} & \num{-12.870} & \num{0.261} \\
Bootstrap por conglomerados (\num{999}) & --- & \num{0.187} & \num{-0.784} & \num{0.893} & \num{-1.102} & \num{-12.456} & \num{0.234} \\
\bottomrule
\end{tabularx}
\par\smallskip{\footnotesize\raggedright $^{\dagger}$: no significativo al 5\,\% en esa especificación.\par}
\end{table}


\figura{sensibilidad}{Coeficientes de las seis explicativas con mayor estadístico $|t|$ bajo
cada especificación, con su intervalo de confianza al 95\,\%. La especificación principal está
resaltada y el bootstrap por conglomerados aparece al final.}

Los resultados son estables. En todas las especificaciones los coeficientes conservan el signo y
un orden de magnitud similar; por ejemplo, el del título universitario (en el Sur) varía entre
\res{rango-grado25}, el de la cobertura sólo pública entre \res{rango-solopublico} y el de la
incidencia entre \res{rango-incidencia}. Las mayores diferencias aparecen con los efectos fijos de
estado y el modelo mixto en la edad mediana y el tamaño del hogar: parte de su asociación con la
mortalidad se produce entre estados, no dentro de ellos. El bootstrap por conglomerados da
intervalos muy parecidos a los paramétricos (por ejemplo, \res{boot-grado25} para el título
universitario frente a \res{ic-grado25}), lo que confirma que la falta de normalidad de los
residuos no afecta a la inferencia.

\section{¿Habría que transformar la respuesta?}

El perfil de verosimilitud de Box-Cox condicionado a las explicativas del modelo final sitúa el
parámetro en $\lambda=\res{boxcox-modelo}$ \res{boxcox-modelo-ic}. El intervalo no contiene
exactamente el 1, pero la transformación óptima es suave y la mejora en la forma de los residuos
no compensa la pérdida de interpretabilidad: los coeficientes dejarían de estar en muertes por
\num{100000} habitantes. Con la inferencia robusta y el tamaño muestral disponibles, se mantiene la
escala original, como pide la orientación.
